\chapter{Modelo del Tiempo de Recorrido}
\label{cap:modelo}

El capítulo anterior sitúa al autobús en el horario: da la hora $\Tpos$ a la que el horario lo pone en el punto donde
está. Falta saber cuánto tardará en llegar desde ahí hasta una parada. Este capítulo explica por qué se trabaja con
diferencias de horario, qué unidad de tiempo se aprende (el tramo entre paradas, por franja), cómo se estima cada
tiempo, qué se hace cuando un tramo no tiene datos y cómo se suman los tramos para obtener el tiempo esperado entre dos
puntos del recorrido.

\section{Por Qué Diferencias de Horario}
\label{sec:diferencias}

Sea $J$ un viaje y $s\in[s_1,s_N]$ un punto de su recorrido, medido en metros desde la cabecera. Se definen tres
funciones:

\begin{itemize}
  \item $T(s)$, la hora a la que el \emph{horario} pone al autobús en $s$. En las paradas, $T(s_i)=a_i$; entre ellas,
        se interpola linealmente (ecuación~\eqref{eq:tpos}).
  \item $R(s)$, la hora \emph{real} a la que el autobús pasa por $s$. Es desconocida salvo en el punto donde se ha
        visto al autobús.
  \item $\Delta(s)=R(s)-T(s)$, el \emph{desfase} en $s$.
\end{itemize}

El autobús se ha visto en $s_p$ a la hora $\tgps=R(s_p)$, y la parada de interés está en $s_d>s_p$, con hora
publicada $a_\dest=T(s_d)$. La cantidad que se quiere predecir es $R(s_d)$.

\begin{proposicion}[El error de las diferencias]
\label{prop:diferencias}
Sean dos estimadores de $R(s_d)$:
\[
\widehat R_{\mathrm{hor}}=T(s_d)\qquad\text{(solo horario)},\qquad
\widehat R_{\mathrm{dif}}=\tgps+\big(T(s_d)-T(s_p)\big)\qquad\text{(diferencias)}.
\]
Sus errores son
\[
R(s_d)-\widehat R_{\mathrm{hor}}=\Delta(s_d),\qquad
R(s_d)-\widehat R_{\mathrm{dif}}=\Delta(s_d)-\Delta(s_p).
\]
Más en general, si $\widehat E$ es cualquier estimación del tiempo real entre $s_p$ y $s_d$ y
$\widehat R=\tgps+\widehat E$, el error es $\big(R(s_d)-R(s_p)\big)-\widehat E$, que no depende de $\Delta(s_p)$.
\end{proposicion}

\begin{proof}
La primera igualdad es la definición de $\Delta$. Para la segunda, como $\tgps=R(s_p)=T(s_p)+\Delta(s_p)$,
\begin{align*}
R(s_d)-\tgps-T(s_d)+T(s_p)&=\big(T(s_d)+\Delta(s_d)\big)-\big(T(s_p)+\Delta(s_p)\big)-T(s_d)+T(s_p)\\
&=\Delta(s_d)-\Delta(s_p).
\end{align*}
La tercera es inmediata: $R(s_d)-\tgps-\widehat E=\big(R(s_d)-R(s_p)\big)-\widehat E$.
\end{proof}

La proposición tiene tres consecuencias que guían todo el diseño.

\begin{enumerate}
  \item \textbf{El retraso acumulado no importa.} Un autobús que va veinte minutos tarde desde la cabecera tiene
        $\Delta(s_p)\approx\Delta(s_d)\approx1200$~s; el horario solo yerra en 20 minutos, las diferencias en casi
        nada. La predicción no necesita saber cuánto retraso lleva el autobús, solo dónde está.
  \item \textbf{El error de forma del horario sí importa.} Como el horario intermedio es una interpolación a velocidad
        constante (sección «\nameref{sec:interpolado}»), $T(s_d)-T(s_p)$ puede estar muy lejos del tiempo real entre esos dos
        puntos: en un tramo de autopista, el doble. Las diferencias de horario eliminan el desfase acumulado, pero no
        este error. Por eso se sustituye $T(s_d)-T(s_p)$ por una estimación aprendida $\widehat E$ del tiempo real
        (el resto de este capítulo).
  \item \textbf{El error crece con la distancia.} El error $\big(R(s_d)-R(s_p)\big)-\widehat E$ suma la variación de
        todos los tramos entre el autobús y la parada. Si los tiempos de los tramos fueran independientes, su varianza
        sería la suma de las varianzas de cada tramo y crecería aproximadamente de forma lineal con el número de
        tramos. Es lo que se observa: la precisión cae de forma regular con el horizonte (capítulo «\nameref{cap:evaluacion}»).
\end{enumerate}

\section{La Unidad de Aprendizaje: El Tramo}
\label{sec:unidad-tramo}

El error de forma del horario se produce tramo a tramo: en cada tramo de autopista el horario se queda corto y en
cada tramo urbano congestionado se queda largo. Si se aprende el tiempo real de cada tramo, el error de forma
desaparece al sumar. Esta es la razón de aprender \emph{por tramo} y no por línea:

\begin{itemize}
  \item un factor único por línea (por ejemplo, «la línea 3414 tarda un 4\,\% menos que su horario») no puede corregir
        un viaje que combina ciudad y autopista, porque la corrección tiene signos distintos en cada parte;
  \item el tramo se identifica por sus dos paradas (definición de tramo, sección «\nameref{sec:notacion}»), de modo que las
        medidas de todas las líneas que lo recorren se suman y el modelo aprende más deprisa;
  \item el tiempo de un tramo es una magnitud física estable (depende de la calle, los semáforos y la demanda de esa
        parada) y fácil de auditar: se puede comparar con lo que tarda uno mismo en recorrerlo.
\end{itemize}

El tiempo de un tramo cambia con la hora y el día. Se aprende por \emph{franja} (definición~\ref{def:franja}): doce
combinaciones de tipo de día (laborable, fin de semana) y periodo del día (madrugada, punta de la mañana, mañana,
mediodía, punta de la tarde, noche). Las fronteras de las franjas siguen los periodos de demanda habituales en el
transporte urbano. Con los datos de octubre de 2026, los tiempos de un mismo tramo difieren entre franjas un 8--13\,\%
de media, y la noche es un 12\,\% más rápida que la punta de la tarde.

\section{Estimador del Tiempo de un Tramo}
\label{sec:estimador-tramo}

\subsection{Las Muestras}

El capítulo «\nameref{cap:aprendizaje}» explica cómo se reconstruye, a partir de las capturas del feed, el instante real en
que un autobús pasa por cada parada. Con dos pasos consecutivos se obtiene una \emph{muestra de tramo}:

\begin{definicion}[Muestra de tramo]
Una muestra de tramo es una tupla $x=(\kappa,\ j,\ b,\ y)$. El primer elemento, $\kappa=u_i{>}u_{i+1}$, es el
identificador del tramo. $j$ es la \emph{pista}: el par vehículo y viaje en una fecha concreta, que identifica un
recorrido físico de un autobús. $b$ es la franja del instante en que el autobús pasó por $p_i$, e $y>0$ es el tiempo
real, en segundos, que tardó en ir de $p_i$ a $p_{i+1}$.
\end{definicion}

\subsection{Un Valor por Pista}

Para cada tramo $\kappa$ y cada \emph{grupo} $g$ (una franja concreta, o el grupo especial \texttt{all} que reúne
todas las franjas) se forma el conjunto de pistas que tienen alguna muestra de $\kappa$ en $g$. Cada pista aporta
\emph{un solo valor}: si una pista tuviera varias muestras del mismo tramo y grupo, se queda la \emph{última} en el
orden del fichero de muestras.\footnote{En la práctica una pista recorre cada tramo una sola vez. La regla de «la
última» importa solo para reproducir exactamente el resultado con datos duplicados; los vectores de prueba la
comprueban (Apéndice~\ref{ap:vectores}).} Sea $Y_{\kappa,g}$ la lista de esos valores y $n_{\kappa,g}=|Y_{\kappa,g}|$
el número de pistas distintas.

Un valor por pista evita que un autobús visto muchas veces pese más que otro visto una vez, y es el sentido en que el
estimador es \emph{robusto}: un autobús que se queda parado veinte minutos aporta un valor extremo, pero uno solo.

\subsection{Mediana y Contracción Hacia el Tramo a Cualquier Hora}
\label{sec:contraccion-tramo}

Para cada tramo $\kappa$ se calcula la mediana baja a cualquier hora,
\begin{equation}
m_{\kappa}=\mediana\big(Y_{\kappa,\texttt{all}}\big),
\end{equation}
y, para cada franja $b$ con $n_{\kappa,b}\ge2$ pistas, la mediana de la franja $m_{\kappa,b}=\mediana(Y_{\kappa,b})$
contraída hacia $m_\kappa$:
\begin{equation}
\label{eq:contraccion-tramo}
\boxed{\;\tau_{\kappa,b}=\redondeo\!\left(\frac{n_{\kappa,b}\,m_{\kappa,b}+n_0\,m_\kappa}{n_{\kappa,b}+n_0}\right),
\qquad n_0=5\;}
\end{equation}
(\texttt{SEGMENT\_BAND\_PRIOR\_TRACKS}). El valor a cualquier hora se guarda como
$\tau_{\kappa,\texttt{all}}=\redondeo(m_\kappa)$, pero solo si el tramo tiene muestras en al menos dos franjas
distintas. Las reglas completas de qué se guarda son:

\begin{enumerate}
  \item Un tramo se considera solo si $n_{\kappa,\texttt{all}}\ge2$ (\texttt{MIN\_TRACKS\_PER\_SEGMENT}).
  \item Una franja $b$ se guarda solo si $n_{\kappa,b}\ge2$.
  \item El grupo \texttt{all} se guarda solo si hay muestras de al menos dos franjas distintas, aunque alguna de ellas
        tenga una sola pista. Con una sola franja, el valor a cualquier hora sería el de esa franja con otro nombre, y
        se usaría en otras franjas para las que no hay ninguna medida: lo medido a mediodía acabaría usándose en hora
        punta.
  \item El tramo se guarda si se guardó al menos un grupo. Se guarda también $n_{\kappa,\texttt{all}}$ como
        información (clave \texttt{n}).
\end{enumerate}

Todos los valores se guardan como enteros (segundos, con $\redondeo$).

La ecuación~\eqref{eq:contraccion-tramo} es una media ponderada entre lo que dice la
franja y lo que dice el tramo a cualquier hora, con pesos $n_{\kappa,b}$ y $n_0$. Es exactamente la forma de la media
a posteriori de un modelo normal jerárquico \citep{gelman2013}: si las pistas de la franja $b$ se distribuyen
alrededor de su media $\mu_b$ con varianza $\sigma^2$, y las medias de las franjas se distribuyen alrededor de la media
del tramo $\mu$ con varianza $\sigma_0^2$, la media a posteriori de $\mu_b$ es
\[
\mathbb{E}[\mu_b\mid\text{datos}]=\frac{n\,\bar y_b+(\sigma^2/\sigma_0^2)\,\mu}{n+\sigma^2/\sigma_0^2}.
\]
La constante $n_0$ representa el cociente $\sigma^2/\sigma_0^2$: cuántas pistas de la franja hacen falta para que su
propia mediana pese lo mismo que la del tramo. Con pocas pistas manda el tramo; con muchas, la franja. El método
sustituye las medias por medianas para conservar la robustez, y fija $n_0$ por validación (no por una estimación de las
varianzas), lo que es más sencillo y se comprobó suficiente.

\begin{proposicion}[Por qué contraer reduce el error]
\label{prop:contraccion}
Sea $\hat\theta_b$ un estimador insesgado de $\theta_b$ con varianza $\sigma^2/n$, y $\mu$ un valor fijo con
$\theta_b-\mu=\beta$. Para el estimador contraído $\tilde\theta_b=w\hat\theta_b+(1-w)\mu$, con $w=n/(n+n_0)$, el error
cuadrático medio es
\[
\mathrm{ECM}(\tilde\theta_b)=w^2\frac{\sigma^2}{n}+(1-w)^2\beta^2 ,
\]
que es menor que el de $\hat\theta_b$ ($\sigma^2/n$) siempre que $\beta^2<\sigma^2\,(2n+n_0)/(n\,n_0)$. En particular,
la ganancia es mayor cuanto menor es $n$.
\end{proposicion}

\begin{proof}
$\tilde\theta_b-\theta_b=w(\hat\theta_b-\theta_b)-(1-w)\beta$; el primer término tiene media cero y varianza
$w^2\sigma^2/n$ y el segundo es constante, de donde sale la expresión. La desigualdad
$w^2\sigma^2/n+(1-w)^2\beta^2<\sigma^2/n$ equivale a $(1-w)^2\beta^2<(1-w^2)\sigma^2/n$, es decir,
$\beta^2<\frac{1+w}{1-w}\cdot\frac{\sigma^2}{n}$, y con $w=n/(n+n_0)$ se tiene $\frac{1+w}{1-w}=\frac{2n+n_0}{n_0}$.
\end{proof}

La proposición~\ref{prop:contraccion} explica el comportamiento medido (sección «\nameref{sec:eval-contraccion}»): la contracción mejora las franjas
con pocas pistas, donde la mediana es ruidosa, y apenas cambia las que tienen muchas. Un $n_0$ demasiado grande sesga
las franjas que de verdad son distintas (el término $\beta^2$), y por eso el valor 10 resultó peor que el 5 a 12--20
minutos. La proposición supone un estimador insesgado y un $\mu$ fijo; el método usa medianas bajas y estima $\mu$
($m_\kappa$) con datos que incluyen las pistas de la propia franja, así que es una justificación heurística, no una
garantía: el valor de $n_0$ lo decide la medida.

\begin{ejemplo}
\label{ej:contraccion}
Un tramo $A{>}B$ tiene dos pistas en \texttt{laborable.punta\_manana} (90 y 100~s) y seis en
\texttt{laborable.mediodia} (60, 58, 62, 64, 59 y 61~s). Las ocho juntas, ordenadas, son 58, 59, 60, 61, 62, 64, 90,
100, y su mediana baja es el elemento $\lfloor 7/2\rfloor=3$: $m_{AB}=61$. Entonces
\begin{align*}
\tau_{AB,\texttt{punta\_manana}}&=\redondeo\!\left(\tfrac{2\cdot90+5\cdot61}{7}\right)=\redondeo(69{,}29)=69,\\
\tau_{AB,\texttt{mediodia}}&=\redondeo\!\left(\tfrac{6\cdot60+5\cdot61}{11}\right)=\redondeo(60{,}45)=60,
\end{align*}
y $\tau_{AB,\texttt{all}}=61$ porque hay dos franjas. Sin contracción, la punta de la mañana valdría 90~s, apoyada solo
en dos autobuses. Un tramo $C{>}D$ con tres pistas, todas de noche (200, 210 y 190~s), guarda solo
\texttt{laborable.noche} $=200$ (la contracción hacia sí mismo no cambia nada) y ningún valor \texttt{all}.
\end{ejemplo}

\subsection{Consulta}

Dado un tramo $\kappa$ y una franja $b$, el tiempo aprendido es
\begin{equation}
\label{eq:consulta-tramo}
\tau(\kappa,b)=\begin{cases}
\tau_{\kappa,b} & \text{si existe},\\
\tau_{\kappa,\texttt{all}} & \text{si no, y si existe},\\
\text{indefinido} & \text{en otro caso.}
\end{cases}
\end{equation}

Se dice que el tramo está \emph{aprendido} en la franja $b$ si $\tau(\kappa,b)$ está definido. Obsérvese que un tramo
con datos solo de noche está aprendido de noche y no lo está a mediodía: entonces se usa el respaldo de la sección
siguiente, no el valor nocturno.

\section{Respaldo por Longitud del Tramo}
\label{sec:respaldo-largo}

Un tramo sin tiempo aprendido no se puede dejar sin estimar, porque la llegada suma todos los tramos del camino. La
primera idea fue usar el horario multiplicado por un factor de ritmo de la línea (sección «\nameref{sec:k}»), pero eso
reproduce el error de forma del horario interpolado. La observación que lleva al respaldo es que ese error depende
sobre todo de la \emph{longitud} del tramo: los tramos cortos son urbanos (el autobús tarda lo que dice el horario o
algo más), y los largos son interurbanos y de autopista (el autobús tarda bastante menos de lo programado, porque la
interpolación reparte el tiempo a velocidad constante).

\subsection{Cubos de Longitud}

La longitud de un tramo es $L=s_{i+1}-s_i$, a partir de \texttt{shape\_dist\_traveled}. Se clasifica en seis cubos
(\texttt{LENGTH\_CUTS\_METERS}):
\begin{equation}
\label{eq:cubo}
c(L)=\begin{cases}
0 & L<300\\
1 & 300\le L<600\\
2 & 600\le L<1000\\
3 & 1000\le L<1500\\
4 & 1500\le L<3000\\
5 & L\ge3000
\end{cases}
\end{equation}
Los cortes son cerrados por abajo: un tramo de exactamente 300~m está en el cubo 1.

\subsection{Proporción Real/Horario por Cubo y Franja}

Para cada muestra de tramo con duración programada $h=a_{i+1}-a_i>0$ y longitud conocida $L>0$, se calcula la
proporción $q=y/h$. Para cada cubo $c$ y grupo $g$ (franja o \texttt{all}) se toma la mediana baja de \emph{todas} las
proporciones del grupo (aquí sin reducir a un valor por pista, porque cada proporción ya es adimensional y compara
tramos distintos):
\begin{equation}
\label{eq:proporcion}
r_{c,g}=\redondeo_3\big(\mediana\{q\}\big),
\end{equation}
siempre que el grupo tenga al menos 20 proporciones (\texttt{MIN\_SAMPLES\_PER\_LENGTH\_RATIO}); si no, $r_{c,g}$ no
existe. Aquí $\redondeo_3(x)=\redondeo(1000x)/1000$ redondea a tres decimales con la misma regla de las mitades. La longitud
$L$ de un tramo, al calibrar, es la mediana baja de las longitudes $s_{i+1}-s_i$ de ese par de paradas en todos los
viajes del horario que lo recorren (puede variar algo entre viajes si los trazados difieren).

\subsection{Consulta}

Al predecir, para un tramo no aprendido con duración programada $h$ y longitud $L$ del viaje concreto,
\begin{equation}
\label{eq:respaldo}
\tau^{\mathrm{L}}(h,L,b)=\clamp\big(\hat r,\ 0{,}3,\ 1{,}5\big)\cdot h,\qquad
\hat r=\begin{cases}
r_{c(L),b} & \text{si existe},\\
r_{c(L),\texttt{all}} & \text{si no.}
\end{cases}
\end{equation}
Si no existe ninguna de las dos proporciones, si $L$ no se conoce o si $L\le0$, $\tau^{\mathrm{L}}$ está indefinido.
Las cotas $[0{,}3;\,1{,}5]$ impiden que un dato corrupto (una proporción de 0,1 o de 2) se traslade a la predicción.
Si el viaje no tiene distancias (sin \texttt{shape\_dist\_traveled}), $L$ no se conoce y este respaldo no se aplica.

La Tabla~\ref{tab:proporciones} muestra las proporciones calibradas el 31 de octubre de 2026. En la columna
\texttt{all} la tendencia es clara y monótona: los tramos de menos de 300~m se recorren en el tiempo programado y los
de más de 3~km en algo más de la mitad. En las franjas, los cubos 3 y 4 quedan casi iguales.

\begin{table}[ht]
\centering
\caption{Proporción Entre Tiempo Real y Tiempo Programado por Cubo de Longitud}\label{tab:proporciones}
\begin{tabular}{@{}clcccc@{}}
\toprule
Cubo & Longitud & $r_{c,\texttt{all}}$ & \shortstack{Laborable\\punta mañana} & \shortstack{Laborable\\mediodía} & \shortstack{Fin de semana\\mañana}\\
\midrule
0 & $<300$~m        & 0,987 & 1,016 & 1,052 & 0,894\\
1 & 300--600~m      & 0,907 & 0,937 & 0,949 & 0,823\\
2 & 600--1\,000~m   & 0,769 & 0,811 & 0,793 & 0,716\\
3 & 1\,000--1\,500~m & 0,647 & 0,665 & 0,659 & 0,606\\
4 & 1\,500--3\,000~m & 0,640 & 0,668 & 0,669 & 0,600\\
5 & $\ge3\,000$~m    & 0,534 & 0,565 & 0,524 & 0,523\\
\bottomrule
\end{tabular}
\nota{Columna $r_{c,\texttt{all}}$: grupo \texttt{all}; las demás, la franja de su título. Muestras acumuladas a 31
de octubre de 2026.}
\end{table}

\section{Factor de Ritmo}
\label{sec:k}

El factor de ritmo $k$ es el último recurso de la cascada y, además, el único modelo cuando no existe ninguna tabla de
tiempos por tramo (por ejemplo, al poner el método en marcha en una red nueva). Entonces las capas 1 y 2 usan
$E=k\,(t_1-t_0)$, la capa 3 no se aplica (las filas sin datos del autobús salen con el horario publicado) y el autobús
que espera en la cabecera se calcula con la ecuación~\eqref{eq:eta-capa1}. Mide cuánto tarda de verdad un autobús
frente a lo que dice el horario, de forma global por franja y, con datos suficientes, por línea y franja.

\subsection{Consulta}

\begin{equation}
\label{eq:k}
k(\ell,D,t)=\clamp\big(\hat k,\ 0{,}6,\ 1{,}2\big),\qquad
\hat k=\begin{cases}
k_{\ell,b} & \text{si existe},\\
k_b & \text{si no, y si existe},\\
1 & \text{en otro caso,}
\end{cases}\qquad b=\franja(D,t),
\end{equation}
donde $\ell$ es la línea del viaje (\texttt{MIN\_FACTOR}, \texttt{MAX\_FACTOR}). El número de línea se escribe
siempre como en el horario, sin ceros a la izquierda: el feed pone \texttt{A0651} y se lee como \texttt{651}, y una
tabla guardada con la otra forma vale igual (\texttt{PaceFactors::lineKey}).

\subsection{Estimación}

Las \emph{muestras de ritmo} (capítulo «\nameref{cap:aprendizaje}») son tuplas $(j,\ \ell,\ v,\ b,\ H,\ Y)$: la pista $j$,
la línea $\ell$, el vehículo $v$, la franja $b$ del instante de la observación, el tiempo que según el horario faltaba
para llegar a una parada, $H=a_\dest-\Tpos\ge120$~s, y el tiempo que de verdad faltaba, $Y=t_{\mathrm{real}}-\tgps>0$.

Para un conjunto de pistas $\mathcal J$ (las de una franja, o las de una línea en una franja), se calcula para cada
pista la proporción de sumas
\begin{equation}
\rho_j=\frac{\sum Y}{\sum H}\qquad\text{(sumas sobre las muestras de la pista $j$)},
\end{equation}
que pondera cada muestra por su duración (una predicción a 15 minutos pesa más que una a 2), y después
\begin{equation}
\label{eq:k-contraido}
\boxed{\;k=\redondeo_3\!\left(\clamp\!\Big(1+\big(\tilde\rho-1\big)\frac{m}{m+10},\ 0{,}6,\ 1{,}2\Big)\right),\qquad
\tilde\rho=\mediana_j\,\rho_j,\quad m=|\mathcal J|\;}
\end{equation}
(\texttt{SHRINKAGE\_TRACKS}). Es una contracción hacia $k=1$ (el horario tal cual) con peso equivalente a diez pistas:
la misma idea que la sección «\nameref{sec:contraccion-tramo}», con el horario como valor común.

Se estima $k_b$ si la franja tiene al menos $m\ge8$ pistas (\texttt{MIN\_TRACKS\_PER\_BAND}), y $k_{\ell,b}$ si la
línea en esa franja tiene al menos 25 pistas y 5 vehículos distintos (\texttt{MIN\_TRACKS\_PER\_LINE},
\texttt{MIN\_VEHICLES\_PER\_LINE}). La exigencia de vehículos distintos evita que el factor de una línea refleje a un
solo conductor o a un autobús averiado.

\begin{ejemplo}
Ocho pistas con proporciones 0,80, 0,82, \dots, 0,94 (en pasos de 0,02) tienen mediana baja $\tilde\rho=0{,}86$ (el
elemento $\lfloor3{,}5\rfloor=3$). Entonces $k=1-0{,}14\cdot8/18=0{,}93778$, que se redondea a $0{,}938$. Con 40
pistas de proporción 0,85, $k=1-0{,}15\cdot40/50=0{,}88$. Con diez pistas de proporciones 1,6, 1,7, \dots, 2,5, la mediana baja es $2{,}0$ y la fórmula
da $1+1{,}0\cdot10/20=1{,}5$, que queda en la cota: $k=1{,}2$.
\end{ejemplo}

Con los datos de octubre de 2026 los factores globales van de $0{,}756$ (fin de semana, punta de la mañana) a $0{,}928$
(laborable, punta de la tarde), y 54 líneas tienen factor propio en alguna franja.

\section{Tiempo Esperado Entre Dos Instantes del Horario}
\label{sec:tiempo-esperado}

Con los tres niveles anteriores se define la función central del método: el tiempo real que se espera que tarde el
autobús en ir desde el punto donde el horario lo pone a la hora $t_0$ hasta el punto donde lo pone a la hora $t_1$.
Los dos extremos se dan como \emph{horas del horario}, no como distancias, porque es lo que proporciona la
localización ($\Tpos$) y lo que define a cada parada ($a_\dest$).

\begin{definicion}[Tiempo de un tramo]
Para el tramo $i$ de $J$, con duración programada $h_i=a_{i+1}-a_i$, longitud $L_i=s_{i+1}-s_i$ (si se conoce),
identificador $\kappa_i$, franja $b$ y factor $k$:
\begin{equation}
\label{eq:cascada}
\tau_i=\begin{cases}
\tau(\kappa_i,b) & \text{si está definido (tiempo aprendido)},\\
\tau^{\mathrm{L}}(h_i,L_i,b) & \text{si no, y si está definido (respaldo por longitud)},\\
k\cdot h_i & \text{en otro caso (factor de ritmo).}
\end{cases}
\end{equation}
\end{definicion}

\begin{definicion}[Tiempo esperado]
\label{def:tiempo-esperado}
Sea $J$ un viaje con llegadas $a_1\le\dots\le a_N$, y sean $t_0<t_1$ dos horas del horario, una franja $b$ y un
factor $k$. Si no existe ninguna tabla de tiempos por tramo (la tabla \texttt{segments} está vacía),
\[
E(J,t_0,t_1,b,k)=k\,(t_1-t_0).
\]
En otro caso,
\begin{equation}
\label{eq:tiempo-esperado}
\boxed{\;E(J,t_0,t_1,b,k)=\sum_{i\in I(t_0,t_1)}\max(0,w_i)\,\tau_i\;}
\end{equation}
con
\begin{equation}
\label{eq:pesos}
w_i=\begin{cases}\dfrac{\min(a_{i+1},t_1)-\max(a_i,t_0)}{a_{i+1}-a_i} & \text{si } a_{i+1}>a_i,\\[2mm] 1 & \text{si } a_{i+1}=a_i,\end{cases}
\end{equation}
donde el conjunto de tramos $I(t_0,t_1)$ se obtiene recorriendo $i=1,\dots,N-1$ en orden y aplicando dos reglas:
se \emph{salta} el tramo si $a_{i+1}\le t_0$ (ya está detrás) y se \emph{termina} el recorrido en el primer tramo con
$a_i\ge t_1$ (ya está después del destino).
\end{definicion}

El peso $w_i$ es la fracción de la duración programada del tramo $i$ que queda entre $t_0$ y $t_1$: vale 1 en los
tramos completos y una fracción en el tramo donde está el autobús. Equivale a suponer que, \emph{dentro de un tramo},
el autobús avanza al ritmo del horario; entre tramos, en cambio, el ritmo es el aprendido. Se probó una curva no lineal
dentro del tramo y no mejoró (sección «\nameref{sec:descartadas}»).

Varias propiedades de la definición son importantes para replicarla exactamente:

\begin{observacion}[Se recorre por horas, no por índices]
Las reglas de salto y de fin se evalúan sobre las horas $a_i$, no sobre los índices de parada. Si dos paradas
consecutivas $p_j$ y $p_{j+1}$ tienen la misma hora programada, el destino $t_1=a_{j+1}=a_j$ termina el recorrido en el
tramo $j$ (porque $a_j\ge t_1$), de modo que ambas paradas reciben la misma llegada estimada.
\end{observacion}

\begin{observacion}[Tramos de duración nula]
Un tramo con $a_{i+1}=a_i$ (frecuente en el GTFS, donde dos paradas próximas comparten minuto) se incluye con peso 1,
y aporta su tiempo aprendido completo, solo si $t_0<a_i$ y $a_i<t_1$. Por eso el horario corregido de la capa 3
(sección «\nameref{sec:capa3}») empieza en $t_0=d_1-1$ y no en $d_1$: así se incluyen los tramos de duración nula del
principio del viaje. Por la misma razón, $E$ no es exactamente aditiva: si $t_0<t_m<t_1$ y hay un tramo de duración
nula con $a_i=t_m$, $E(t_0,t_1)$ lo cuenta y $E(t_0,t_m)+E(t_m,t_1)$ no.
\end{observacion}

\begin{observacion}[La espera en la parada]
La duración de un tramo, $h_i=a_{i+1}-a_i$, incluye la espera programada en $p_i$, mientras que $\Tpos$ se interpola
entre $d_i$ y $a_{i+1}$ (sección «\nameref{sec:trazado}»). Si $d_i>a_i$, un autobús que acaba de salir de $p_i$
($\Tpos=d_i$) recibe $w_i=(a_{i+1}-d_i)/h_i<1$ y no cuenta la parte del tiempo aprendido que corresponde a esa
espera. En la base de Bizkaibus todas las paradas tienen $d_i=a_i+2$~s, así que el efecto es de unos segundos; los
vectores de prueba, con $a_i=d_i$, no lo ejercitan.
\end{observacion}

\begin{observacion}[Franja única]
Todo el camino se evalúa con una sola franja $b$, la del instante de la observación, aunque el autobús vaya a cruzar a
otra franja antes de llegar. Calcular la franja de cada tramo a la hora prevista de paso por él se probó y no mejoró
(sección «\nameref{sec:descartadas}»).
\end{observacion}

\begin{ejemplo}
\label{ej:tiempo-esperado}
En el viaje sintético del ejemplo~\ref{ej:posicion}, con franja \texttt{laborable.punta\_manana} y $k=0{,}9$, los
tramos aprendidos son $S0{>}S1$ (110~s en esa franja) y $S1{>}S2$ (140~s), y las proporciones del cubo~2 y del cubo~4
en esa franja son $0{,}811$ y $0{,}668$. Desde $t_0=\Tpos=28\,950$ (a mitad del tramo $S1{>}S2$):
\begin{itemize}
  \item hasta $S2$ ($t_1=29\,010$): $w=(29\,010-28\,950)/120=0{,}5$ en $S1{>}S2$, luego $E=0{,}5\cdot140=70$~s;
  \item hasta $S3$ ($t_1=29\,310$): se añade $S2{>}S3$, no aprendido, de 1\,500~m (cubo 4) y 300~s programados:
        $0{,}668\cdot300=200{,}4$~s, luego $E=270{,}4$~s;
  \item hasta $S4$ ($t_1=29\,430$): se añade $S3{>}S4$, de 700~m (cubo 2) y 120~s: $0{,}811\cdot120=97{,}32$~s, luego
        $E=367{,}72$~s.
\end{itemize}
Sin distancias (viaje sin trazado), los dos últimos tramos caen al factor de ritmo: $0{,}9\cdot300=270$ y
$0{,}9\cdot120=108$~s, y $E$ vale $70$, $340$ y $448$~s.
\end{ejemplo}

\section{La Cascada Completa}
\label{sec:cascada}

La Tabla~\ref{tab:cascada} resume qué nivel del modelo se usa en cada caso. Todos los niveles devuelven un tiempo en
segundos para un tramo; el tiempo esperado de un camino es la suma ponderada de la
ecuación~\eqref{eq:tiempo-esperado}, y cada tramo del camino puede salir de un nivel distinto.

\begin{table}[ht]
\centering
\caption{Cascada de Estimación del Tiempo de un Tramo}\label{tab:cascada}
\begin{tabularx}{\textwidth}{@{}c>{\raggedright\arraybackslash}X>{\raggedright\arraybackslash}X>{\raggedright\arraybackslash}X@{}}
\toprule
Nivel & Se usa si & Valor & Aprende de\\
\midrule
1 & El tramo tiene valor en la franja & $\tau_{\kappa,b}$, contraído hacia el tramo & Pasos reales de $\ge2$ pistas en la franja\\
2 & El tramo tiene valor a cualquier hora (y muestras de $\ge2$ franjas) & $\tau_{\kappa,\texttt{all}}$ & Pasos reales de $\ge2$ pistas\\
3 & Se conoce la longitud y hay proporción para su cubo & $\clamp(\hat r,0{,}3,1{,}5)\cdot h$, con $\hat r=r_{c,b}$ o, si no, $r_{c,\texttt{all}}$ & $\ge20$ proporciones del cubo en el grupo\\
4 & Siempre & $k\cdot h$, con $k\in[0{,}6;\,1{,}2]$ & $\ge8$ pistas (franja) o $\ge25$ y 5 vehículos (línea)\\
\bottomrule
\end{tabularx}
\end{table}

El diseño en cascada responde a un compromiso entre sesgo y varianza: cada nivel es menos específico que el anterior
(más sesgo) pero se apoya en más datos (menos varianza). El nivel que se usa es siempre el más específico que tiene
datos suficientes, y los umbrales de «suficiente» se han fijado por validación cruzada (capítulos «\nameref{cap:aprendizaje}»
y «\nameref{cap:evaluacion}»). Con los datos de octubre de 2026, la tabla tiene 2\,454 tramos aprendidos, el 86\,\% de las
predicciones tiene todos los tramos de su camino en los niveles 1 o 2, y, como los seis cubos tienen proporción a cualquier hora, el nivel 4 solo actúa en la práctica
en viajes sin distancias a lo largo del trazado.
